\section{Tổng kết}

Trong project này, nhóm đã xây dựng một hệ thống nhận diện và ẩn danh thông tin cá nhân trong văn bản bằng phương pháp rule-based. Hệ thống tập trung vào việc phát hiện các thực thể quan trọng như \texttt{NAME}, \texttt{EMAIL}, \texttt{PHONE}, \texttt{ADDRESS}, \texttt{USERNAME} và \texttt{URL}. Thay vì sử dụng các mô hình machine learning hoặc deep learning, project áp dụng các kỹ thuật xử lý chuỗi như regular expression, keyword matching, token scanning và xử lý overlap giữa các entity để giải quyết bài toán NER.

Trong quá trình thực hiện, nhóm đã xây dựng được một pipeline xử lý tương đối hoàn chỉnh, bao gồm các bước preprocessing dữ liệu, phát hiện entity bằng regex, phát hiện entity bằng ngữ cảnh, gộp kết quả, xử lý entity chồng lấn, ẩn danh văn bản và đánh giá kết quả bằng các chỉ số Precision, Recall và F1-score. Việc chia hệ thống thành nhiều module độc lập giúp các thành viên có thể phát triển song song, đồng thời hỗ trợ quá trình kiểm thử và mở rộng sau này trở nên dễ dàng hơn.

Kết quả đạt được cho thấy phương pháp rule-based hoạt động tương đối hiệu quả đối với các thực thể có cấu trúc rõ ràng như email, URL và số điện thoại. Các loại thực thể này thường có format tương đối cố định nên phù hợp với phương pháp regular expression. Đối với các thực thể phụ thuộc nhiều vào ngữ cảnh như tên người hoặc địa chỉ, hệ thống vẫn có thể phát hiện được bằng cách sử dụng các luật thủ công dựa trên keyword và context. Ngoài ra, cơ chế xử lý overlap giữa các entity cũng giúp giảm một phần lỗi nhận diện trùng lặp hoặc phát hiện sai trong văn bản.

Bên cạnh những kết quả đạt được, project vẫn còn tồn tại một số hạn chế nhất định. Vì sử dụng rule-based method nên chất lượng nhận diện phụ thuộc rất nhiều vào tập luật được xây dựng thủ công. Nếu văn bản có cấu trúc quá đa dạng hoặc xuất hiện các format hiếm gặp, hệ thống có thể bỏ sót entity hoặc nhận diện sai. Đặc biệt, các thực thể như tên người và địa chỉ thường có rất nhiều biến thể khác nhau nên khó đạt độ chính xác cao chỉ bằng các luật đơn giản. Ngoài ra, phương pháp đánh giá exact match trong bài toán NER cũng khá nghiêm ngặt, chỉ cần lệch một phần nhỏ vị trí span thì prediction sẽ bị xem là sai hoàn toàn.

Mặc dù vậy, project vẫn thể hiện rõ vai trò quan trọng của bài toán NER trong thời đại AI hiện nay. Khi lượng dữ liệu văn bản ngày càng lớn và được sử dụng trong nhiều hệ thống khác nhau, nhu cầu bảo vệ thông tin cá nhân ngày càng trở nên cần thiết. Các kỹ thuật NER có thể hỗ trợ tự động phát hiện và ẩn danh dữ liệu nhạy cảm trước khi lưu trữ, chia sẻ hoặc xử lý tiếp. Điều này đặc biệt quan trọng trong các lĩnh vực như tài chính, y tế, giáo dục, thương mại điện tử và mạng xã hội, nơi dữ liệu cá nhân thường xuyên được xử lý với số lượng lớn.

Thông qua project này, nhóm không chỉ hiểu rõ hơn về bài toán nhận diện thực thể trong xử lý ngôn ngữ tự nhiên mà còn có cơ hội áp dụng các thuật toán xử lý chuỗi, regular expression và thiết kế pipeline xử lý dữ liệu thực tế. Đây cũng là nền tảng quan trọng để tiếp cận các hướng nghiên cứu hoặc hệ thống NER phức tạp hơn trong tương lai.

\section{Hướng mở rộng trong tương lai}

Trong tương lai, hệ thống có thể được mở rộng theo nhiều hướng khác nhau nhằm cải thiện độ chính xác, khả năng tổng quát hóa và khả năng ứng dụng thực tế của bài toán NER. Một trong những hướng phát triển quan trọng là bổ sung thêm rule và dictionary để tăng khả năng nhận diện tên riêng, địa danh và địa chỉ. Việc sử dụng danh sách first name, last name, tên đường hoặc địa danh phổ biến có thể giúp giảm false positive và tăng độ chính xác cho các context-based detector.

Ngoài ra, project hiện chủ yếu hoạt động trên dữ liệu tiếng Anh nên vẫn chưa tối ưu cho tiếng Việt. Trong tương lai, hệ thống có thể được mở rộng sang dữ liệu tiếng Việt bằng cách xây dựng thêm các rule phù hợp với cấu trúc tên người, số điện thoại và địa chỉ tại Việt Nam. Đây là hướng mở rộng có ý nghĩa thực tế cao vì nhiều hệ thống trong nước hiện vẫn gặp khó khăn trong việc xử lý dữ liệu tiếng Việt.

Một hướng cải tiến khác là tối ưu quá trình xử lý overlap giữa các entity. Hiện tại, project chủ yếu sử dụng cơ chế priority-based greedy selection để xử lý entity chồng lấn. Trong tương lai, có thể áp dụng các kỹ thuật như interval scheduling hoặc weighted interval optimization để lựa chọn tập entity tối ưu hơn. Ngoài ra, fuzzy matching cũng có thể được tích hợp nhằm xử lý các trường hợp entity bị viết tắt hoặc có sai lệch nhỏ về format.

Bên cạnh đó, hệ thống có thể được mở rộng thêm nhiều loại entity mới ngoài phạm vi hiện tại, chẳng hạn như tên tổ chức, mã số định danh, biển số xe, tài khoản ngân hàng hoặc thông tin học tập. Việc mở rộng tập entity sẽ giúp hệ thống phù hợp hơn với các bài toán anonymization thực tế trong doanh nghiệp hoặc các hệ thống lưu trữ dữ liệu lớn.

Một hướng phát triển quan trọng khác là kết hợp rule-based method với machine learning hoặc deep learning nhằm tận dụng ưu điểm của cả hai phương pháp. Rule-based approach có tính giải thích cao và hoạt động tốt với các mẫu cố định, trong khi machine learning có khả năng học được các pattern phức tạp và tổng quát hóa tốt hơn trên dữ liệu đa dạng. Việc kết hợp hybrid approach có thể giúp cải thiện đáng kể hiệu năng của hệ thống NER trong các ngữ cảnh thực tế.

Ngoài phần xử lý backend, project cũng có thể được mở rộng bằng cách xây dựng giao diện web local để trực quan hóa kết quả nhận diện entity. Người dùng có thể nhập văn bản trực tiếp lên website, sau đó hệ thống sẽ highlight các thực thể được phát hiện và hiển thị phiên bản đã được ẩn danh. Điều này giúp project dễ demo hơn và tăng tính ứng dụng thực tế của hệ thống.

Tóm lại, dù vẫn còn nhiều hạn chế, project đã xây dựng được một nền tảng tương đối hoàn chỉnh cho bài toán nhận diện và ẩn danh thông tin cá nhân bằng phương pháp rule-based. Các hướng mở rộng trong tương lai không chỉ giúp cải thiện độ chính xác của hệ thống mà còn mở ra khả năng ứng dụng rộng hơn trong các bài toán NLP và bảo vệ dữ liệu cá nhân trong thời đại AI.